\Needspace{8\baselineskip}\originalchapter{2003 作业三: 漂移、镜与环轨道}
题源: 23 Sep 2003, due 30 Sep. 以下导引中心结果要求回旋尺度小于场变化尺度.
\begin{exercise}\label{ex:03-3-1}
原题1. 粒子在xy面, $\mathbf B=B(x)\mathbf e_z$缓变, $\mathbf E=E\mathbf e_y$, E/B<c. 证明grad-B漂移使E做功, 由能量说明 $mv_\perp^2/(2B)$守恒.
\end{exercise}
\begin{solution}
$\mathbf v_E=(E/B)\mathbf e_x$, $\mathbf v_{\nabla B}=(\mu B'/(qB))\mathbf e_y$, $\mu=mv_\perp^2/(2B)$. 因而 $\dot W=qE v_{\nabla B,y}=\mu EB'/B$, 而随中心 $\dot B=B'v_{E,x}=EB'/B$. 得 $\dot W=\mu\dot B$, $\dd(W/B)/\dd t=0$. 这里$W$与$v_\perp$是相对导引中心的回旋能和速度, 此非相对论演示另采用$mv_E^2/2$为高阶小量的排序. 仅有$E/B<c$不足以保证这一点; 若$v_E$与回旋速度同阶, 还需极化与惯性修正, 不能称实验室总动能除$B$严格守恒.
\end{solution}
\begin{exercise}\label{ex:03-3-2}
原题2. E\ensuremath{\times}B速度不能超c. 当E/B>c, 用 $\dd(m\mathbf v/\sqrt{1-v^2/c^2})/\dd t=q(\mathbf E+\mathbf v\times\mathbf B)$ 或其他方法说明轨道, 并画B沿z、E沿x、无z运动的形状.
\end{exercise}
\begin{solution}
取沿 $\mathbf E\times\mathbf B$ 的Lorentz boost $\mathbf u=-c^2B/E\,\mathbf e_y$, |u|<c, 则新系B\ensuremath{{}^{\prime}}=0、$E'=\sqrt{E^2-c^2B^2}$. 这说明运动是受电场持续加速而非闭合回旋. 新系 $p_x'=p_{x0}'+qE't'$, $p_y'$常数. 以p\ensuremath{{}^{\prime}}x0=0为例, 令 $A=\sqrt{m^2c^2+p_y'^2}$,
\begin{equation}
 x'-x'_0=\frac c{qE'}[\sqrt{A^2+(qE't')^2}-A],\quad
 y'-y'_0=\frac{cp_y'}{qE'}\operatorname{arsinh}\frac{qE't'}A.
\end{equation}
逆boost得 $y=\gamma_u(y'+ut')$, $t=\gamma_u(t'+uy'/c^2)$, x=x\ensuremath{{}^{\prime}}. 长时 $v_x\to\operatorname{sgn}(qE)c/\gamma_u$, $v_y\to u$, 总速c. 初始条件不同可改变短时弯曲, 均不能采用E/B的超光速漂移.
\begin{center}\begin{tikzpicture}[scale=.7]\draw[->](0,0)--(5,0)node[right]{$x$};\draw[->](0,0)--(0,-3)node[below]{$y$};\draw[thick,structurecolor](0,0)..controls(.7,0)and(1.4,-.3)..(2.4,-1.2)..controls(3.3,-1.9)and(4,-2.4)..(4.6,-2.8);\node at(3.1,.6){$qE>0$, 渐近直线};\end{tikzpicture}\end{center}
图为电场主导的定性轨道, 精确曲线由上述初值和boost确定.
\end{solution}
\begin{exercise}\label{ex:03-3-3}
原题3. 镜中心B0、喉B1, 求被俘获电子比例, f0分别为 (a) Maxwell; (b) $[1+v^2/v_t^2]^{-2}$; (c) $(v_\perp/v)\exp[-mv^2/(2T)]$.
\end{exercise}
\begin{solution}
能量和\ensuremath{\mu}守恒给俘获条件 $|\cos\alpha|<u_0=\sqrt{1-B_0/B_1}$. (a)(b)均各向同性, 角均匀du, 比例u0, 与径向速度函数无关. (c)角权重 $\sqrt{1-u^2}$, 除以 $\int_{-1}^1\sqrt{1-u^2}\dd u=\pi/2$, 得 $[2/\pi][\arcsin u_0+u_0\sqrt{1-u_0^2}]$. 这只按给定分布计数, 不处理碰撞补充损失锥.
\end{solution}
\begin{exercise}\label{ex:03-3-4}
原题4. 均匀B沿z, 轴线上线电荷Q每单位长度. (a)导引中心轨迹及xy返回周期; (b)B以小常数Bdot变化, 求中心半径rg和动能演化; (c)动能加q\ensuremath{\phi}是否守恒, q、m是否影响?
\end{exercise}
\begin{solution}
$E_r=Q/(2\pi\epsilon_0r)$, $v_\theta=-Q/(2\pi\epsilon_0Br)$, 固定r圆轨道, 周期 $4\pi^2\epsilon_0Br^2/|Q|$. 慢变B有Faraday $E_\theta=-r\dot B/2$, 故 $\dot r=-r\dot B/(2B)$, $r^2B=r_i^2B_i$, $r=r_i\sqrt{B_i/B}$. 绝热回旋能 $W_\perp=\mu B$, 平行动能保持. 一阶导引中心忽略二阶漂移动能; 如需它, 应同时保留同阶轨道修正. 电势 $\phi=-Q\ln(r/r_{\rm ref})/(2\pi\epsilon_0)$, 所以
$\dd(W+q\phi)/\dd t=\mu\dot B+qQ\dot B/(4\pi\epsilon_0B)$, 一般不为零, 能量来自感应电场的外电源. 中心几何漂移到此阶与q/m无关, 但回旋频率、\ensuremath{\mu}及电势能依粒子参数, 不能说所有能量变化都无关.
\end{solution}
\begin{exercise}\label{ex:03-3-5}
原题5. 无限长均匀氢等离子体, Te=Ti=T, 环绕载流I无限直圆柱导体. 忽略等离子体对B的扰动, 求 (a)平均电子漂移; (b)离子漂移; (c)电流密度; (d)加分: 若不忽略其电流, 内边界r=a, 求B(r).
\end{exercise}
\begin{solution}
此题``平均漂移''有两种计量, 必须区分. 对导引中心, $B_\theta=\mu_0I/(2\pi r)$, gradB和曲率同向组合, Maxwell平均给 $\langle v_{gc,z}\rangle=2T/(qBr)$; 电子负、离子正, $J_{gc,z}=4nT/(Br)$. 但均匀\textbf{粒子位置}Maxwell $f(\mathbf r,\mathbf v)=f(v^2)$ 是精确静态Vlasov解, 实际 $\int q\mathbf vf\dd^3v=0$. 磁化电流抵消上述中心电流: $\mathbf M=-p\mathbf B/B^2$, 总p=2nT, $\nabla\times\mathbf M=-4nT/(Br)\mathbf e_z$. 故自洽均匀压力的体电流仍零, 外导体场仍1/r.

若只把Jgc错误当全部电流再解安培, 得 $B^2=4\mu_0nT+C/r^2$; 这个形式不满足均匀p的全流体力平衡, 因而不能当作(d)的无条件物理解. 若题意另指定均匀\textbf{中心}分布而非实际粒子分布, 要给其坐标测度和磁化/边界条件才能求完整自洽场. 本解给出中心估计和均匀粒子体系的正确结论, 明示原题省略的区别.
\end{solution}
